% =====================================================================
%  Betriebsanweisung (Gefahrstoff) Phrozen Water-Washable Rapid Black
%  (wird von SLA_Drucker_Einweisung.tex eingebunden)
%  Grundlage: Kennzeichnung von Phrozen Water-Washable Resin Schwarz
%  (Rapid Black) laut EU-Händlerangabe (3djake.de, Stand 28.09.2026):
%  GHS05, GHS07, GHS08, GHS09, Signalwort Gefahr,
%  H315, H317, H318, H335, H360FD, H373, H410.
%  Das Sicherheitsdatenblatt des Herstellers muss im FabLab vorliegen.
% =====================================================================

\BAgefahrstofffarbe
\renewcommand{\BAgrundlage}{gemäß § 14 GefStoffV und TRGS 555}
\renewcommand{\BAgeraetlabel}{Gefahrstoff:}
\renewcommand{\BAgeraet}{Phrozen Water-Washable Rapid Black -- flüssiges Photopolymer-Harz
  auf Acrylatbasis}
\renewcommand{\BAnummer}{BA-GS-01}
% Signalwort im Kopf, Piktogramme im Abschnitt Gefahren
\renewcommand{\BAgeraetbilder}{{\bfseries\large GEFAHR}}

\BAkopf{Drucken, Waschen im Form Wash (Wasser), Nachhärten, Harz einfüllen, Entsorgung}

% ---------- 1 Gefahren -----------------------------------------------
\BAblock{GEFAHREN FÜR MENSCH UND UMWELT}{GHS08,GHS05,GHS07,GHS09}{%
  \begin{itemize}
    \item \textbf{Kann die Fruchtbarkeit beeinträchtigen. Kann das Kind im
          Mutterleib schädigen} (H360FD).
    \item \textbf{Verursacht schwere Augenschäden} (H318).
    \item \textbf{Kann allergische Hautreaktionen verursachen} (H317). Eine
          Allergie gegen Acrylate bleibt ein Leben lang bestehen.
    \item Verursacht Hautreizungen (H315). Kann die Atemwege reizen (H335).
          Kann die Organe schädigen bei längerer oder wiederholter
          Exposition (H373).
    \item \textbf{Sehr giftig für Wasserorganismen} mit langfristiger
          Wirkung (H410).
    \item Das gilt auch für \textbf{Waschwasser} und für nicht nachgehärtete
          Druckteile.
  \end{itemize}}

% ---------- 2 Schutzmaßnahmen ----------------------------------------
\BAblock{SCHUTZMASSNAHMEN UND VERHALTENSREGELN}{M004,M009,P022}{%
  \begin{itemize}
    \item \textbf{Schwangere und Stillende dürfen nicht mit dem Harz
          arbeiten} (Mutterschutzgesetz).
    \item Eine \textbf{Schutzbrille} und \textbf{Nitrilhandschuhe} tragen;
          verschmutzte Handschuhe sofort wechseln.
    \item Haut- und Augenkontakt vermeiden, Dämpfe nicht einatmen. Den Raum
          gut lüften, Drucker und Flasche geschlossen halten.
    \item Die Teile erst nach Waschen und Nachhärten ohne Handschuhe anfassen.
    \item Nach der Arbeit die Hände waschen. Essen und Trinken sind verboten.
    \item Das Harz dicht verschlossen, unter Verschluss und dunkel bei
          15--35\,\textcelsius{} lagern.
  \end{itemize}}

% ---------- 3 Gefahrfall ---------------------------------------------
\BAblock{VERHALTEN IM GEFAHRFALL}{W001,F001}{%
  \begin{itemize}
    \item Verschüttetes Harz mit Papiertüchern oder Bindemittel aufnehmen,
          nicht in den Abfluss gelangen lassen und einen Betreuer informieren.
    \item Verschmutzte Kleidung sofort ausziehen.
    \item Bei einem Brand mit Wassersprühstrahl, CO\textsubscript{2} oder
          Pulver löschen. Die Brandgase (CO, NO\textsubscript{x}) nicht einatmen.
          \textbf{Feuerwehr: 112}
  \end{itemize}}

% ---------- 4 Erste Hilfe --------------------------------------------
\BAblock{ERSTE HILFE}{E003,E011}{%
  \begin{itemize}
    \item \textbf{Augen:} sofort mindestens 10 Minuten bei geöffnetem Lid mit
          Wasser spülen, Kontaktlinsen entfernen und \textbf{sofort zum Augenarzt}.
    \item \textbf{Haut:} mit viel Wasser und Seife waschen. Bei Rötung
          oder Ausschlag einen Arzt aufsuchen.
    \item \textbf{Einatmen:} an die frische Luft gehen, bei Beschwerden zum Arzt.
    \item \textbf{Verschlucken:} den Mund ausspülen, kein Erbrechen
          herbeiführen, einen Arzt aufsuchen.
    \item Dem Arzt das Sicherheitsdatenblatt oder die Flasche zeigen.
  \end{itemize}\vspace{1.5mm}
  \textbf{Notruf: 112}\qquad Giftnotruf München: 089\,/\,19240\qquad
  FabLab: 09131\,/\,85\,28013}

% ---------- 5 Entsorgung (mit Fußzeile) --------------------------------
\BAletzterblock{SACHGERECHTE ENTSORGUNG}{}{%
  \begin{itemize}
    \item Flüssiges Harz und \textbf{Waschwasser nie in den Abfluss} gießen.
          Das Waschwasser bleibt im Form Wash.
    \item Harzreste, Tücher und Handschuhe mit UV-Licht vollständig
          aushärten und danach im Restmüll entsorgen.
    \item Altes Waschwasser entsorgen nur die Betreuer.
  \end{itemize}}
